\begin{table}[!h]
\centering
\caption{\label{tab:candidate_women}Candidate Characteristics Among Women: Quota vs. Open Seats}
\centering
\fontsize{8}{10}\selectfont
\begin{threeparttable}
\begin{tabular}[t]{lllllll}
\toprule
\multicolumn{1}{c}{ } & \multicolumn{3}{c}{Women winners} & \multicolumn{3}{c}{Women candidates} \\
\cmidrule(l{3pt}r{3pt}){2-4} \cmidrule(l{3pt}r{3pt}){5-7}
Characteristic & Open & Quota & Difference & Open & Quota & Difference\\
\midrule
Age & 42.8 & 41.3 & 1.6 & 40.6 & 39.9 & 0.7\\
Graduate (\%) & 12.2 & 12.7 & -0.5 & 17.3 & 13.3 & 4.0\\
Unemployed (\%) & 53.1 & 61.1 & -8.0 & 54.8 & 61.6 & -6.8\\
Log assets & 12.6 & 12.6 & 0.0 & 12.1 & 12.1 & 0.0\\
\bottomrule
\end{tabular}
\begin{tablenotes}
\item \textit{Note: } 
\item Descriptive open-minus-quota comparisons, conditional on candidacy or winning. Graduation and unemployment are percentages; their differences are percentage points. Other rows use their stated units. Graduation rates exclude records with missing education. Age means exclude recorded ages outside 21--110 years; other outcomes retain those records. Assets winsorized at the 10th and 90th percentiles of the candidate sample before taking logs. These comparisons do not identify a causal effect on candidate quality.
\end{tablenotes}
\end{threeparttable}
\end{table}
